\chapter{Cosa si conosce davvero, in Veneto}
\label{cap:cosa_si_conosce}

Prima di chiedersi cosa manca, conviene fare ordine su cosa esiste, con la maggiore precisione possibile. Il territorio veneto, inteso qui nel senso geografico che tiene conto dei confini naturali piuttosto che amministrativi, non è affatto privo di testimonianze della frequentazione preistorica: possiede siti documentati in ogni periodo, dall'ultima glaciazione fino all'età del ferro, distribuiti in modo disomogeneo tra montagna e pianura, tra aree intensamente indagate e aree che per ragioni storiche e logistiche non hanno ancora ricevuto un'attenzione sistematica comparabile. Questo capitolo li ripercorre provincia per provincia, risalendo per quanto possibile ai lavori originali, con le cifre, le date e le riserve che quei lavori stessi dichiarano.

Una precisazione metodologica preliminare: la ripartizione per provincia moderna è uno strumento pratico, non una categoria preistorica. Le popolazioni paleolitiche e mesolitiche non seguivano confini amministrativi, e molte delle aree più interessanti si trovano nelle fasce di confine tra province o in territori che la geografia fisica renderebbe meglio descritti in relazione ai bacini fluviali. Il concetto di \emph{Veneto geografico esteso} usato in questo libro include, per le stesse ragioni, alcune zone oggi trentine o lombarde ma storicamente legate allo stesso sistema di valli e passi: la sponda settentrionale del Garda, il Primiero, la Valsugana, Folgaria.

% ============================================================
\section{Verona: le Basse Colline, il Garda, la Valpolicella}
% ============================================================

La provincia di Verona ospita due dei tre siti di arte rupestre o arte preistorica documentati con maggiore solidità nell'intero Veneto, e uno dei siti paleolitici più antichi d'Europa. Per questa ragione, e perché la storia della loro scoperta è essa stessa parte del problema che questo libro discute, vengono trattati con il maggior dettaglio consentito dalle fonti disponibili.

\subsection{La grotta di Fumane: quarantamila anni di arte figurativa}

La grotta di Fumane, in Valpolicella, non lontana da Verona, è stata scavata sotto la direzione di Alberto Broglio dell'Università di Ferrara a partire dagli anni novanta del Novecento. Dai livelli attribuiti alla fase più antica della presenza della nostra specie in Europa sono stati recuperati sei frammenti di roccia caduti dalla volta per crioclastismo, ritrovati capovolti sul terreno, dipinti con ocra rossa. Due frammenti sono i più noti e i più discussi in letteratura: una figura antropomorfa alta circa diciotto centimetri, vista frontalmente, che indossa quello che sembra un copricapo o una maschera con due protuberanze simili a corna, con in mano un oggetto interpretato come votivo; e un animale a quattro zampe visto di profilo, verosimilmente un mustelide o un felide, la cui identificazione precisa resta incerta \citep{broglio2003, broglio2005}.

La datazione radiometrica sui livelli associati si aggira attorno ai quarantunomila anni prima dell'anno zero: tra le più antiche mai ottenute per opera figurativa in Europa, e un dato che colloca Fumane in un piccolo gruppo di siti di riferimento mondiale per lo studio delle origini dell'arte figurativa umana.

Fumane non è, in senso stretto, arte rupestre all'aperto: è arte parietale in ambiente di grotta, un fenomeno diverso per almeno due ragioni. La prima è la condizione di conservazione: una grotta protegge dagli agenti atmosferici che, all'aperto, cancellano progressivamente un'incisione o una pittura. La seconda è il contesto di ritrovamento: Fumane è nota da uno scavo stratigrafico controllato, con datazione diretta dei livelli associati. Va tenuto fermo cosa Fumane dimostra senza possibilità di dubbio: che la pratica di produrre immagini su superfici rocciose era presente sul territorio veneto già quarantamila anni fa. Va altrettanto tenuto fermo cosa Fumane non dimostra: non è una prova dell'esistenza di incisioni all'aperto nelle fasi successive, separate da decine di millenni e da tradizioni culturali completamente diverse.

\subsection{Il Monte Luppia e le incisioni del Benaco}

Il complesso di arte rupestre all'aperto più importante del Veneto si trova sulle pendici collinari che si affacciano sul lago di Garda, in particolare sul Monte Luppia, tra Punta San Vigilio e Malcesine. Le prime rocce incise furono individuate nel 1964 da Mario Pasotti, insegnante, dopo una visita alle incisioni della Valcamonica che gli fece riconoscere segni simili sulle rocce vicino a casa. Il dettaglio vale la pena sottolinearlo per intero: Pasotti non era un archeologo di professione, e la sua scoperta non nasce da una prospezione pianificata a tavolino, ma dal riconoscimento di un pattern già visto altrove, applicato a un territorio che nessuno aveva ancora guardato con quello sguardo. Il lavoro di catalogazione fu proseguito e ampliato da Fabio Gaggia, che nel 1985 costituì un centro di studi dedicato al territorio benacense e ne rimane il principale studioso \citep{gaggia1982, gaggia2002}.

Il catalogo oggi comprende più di duecentocinquanta rocce incise, per un totale di almeno tremila figure: guerrieri, cavalieri, imbarcazioni, armi, simboli solari, coppelle isolate o in gruppo, e un numero consistente di segni geometrici di incerta interpretazione. La cronologia copre un arco lungo: le fasi più antiche vengono attribuite all'età del bronzo (la Pietra di Castelletto e la Pietra delle Griselle, collocate nel secondo millennio prima dell'anno zero), le fasi intermedie all'età del ferro (con nuclei a Senge di Marciaga, nel comune di Costermano, e a Crero, nel comune di Brenzone), e uno strato più recente arriva fino a incisioni chiaramente medievali o moderne. Il museo del castello scaligero di Torri del Benaco colloca esplicitamente questo complesso subito dopo la Valcamonica e il Monte Bego per importanza nel panorama alpino.

Il dato più importante per il tema di questo libro, però, non è la dimensione del catalogo: è che quel catalogo è iniziato nel 1964 e non prima. Le rocce del Monte Luppia erano lì, esposte, raggiungibili, da millenni. Quello che è cambiato in quell'anno non è il territorio, ma la disponibilità di uno sguardo capace di riconoscere un'incisione intenzionale: un caso da manuale del meccanismo del research bias, documentabile con un nome, una data e una motivazione dichiarata dallo stesso scopritore.

\subsection{Riparo Tagliente e la Lessinia}

La Lessinia, altopiano carsico a nord di Verona, ospita una serie di siti paleolitici documentati con scavi sistematici. Il più importante è Riparo Tagliente, nella Valle delle Falci, scavato a partire dagli anni sessanta e poi con continuità nei decenni successivi. Il sito ha restituito una lunga sequenza occupazionale che va dal Musteriano (Neandertal) fino all'Epigravettiano finale, con una presenza documentata attorno ai diciottomila anni prima dell'anno zero per le fasi più recenti. La fauna cacciata include cavallo, cervo, stambecco, camoscio e rinoceronte lanoso per le fasi più antiche. Non è stata trovata arte parietale \emph{in situ}, ma il sito ha restituito elementi di ornamento personale e strumenti ossei lavorati.

Nella stessa area: Riparo Mezzena (sequenza musteriana), Grotta della Ghiacciaia (resti faunistici del Pleistocene superiore), e la Grotta di San Bernardino (occupazione paleolitica con resti di Homo sapiens datati a circa ventottomila anni prima dell'anno zero, tra i più antichi della regione) \citep{bartolomei1992}.

% ============================================================
\section{Vicenza: l'Altopiano di Asiago e i Colli Berici}
% ============================================================

\subsection{La Val d'Assa e le incisioni dell'altopiano}

Sull'altopiano di Asiago, presso Canove e Roana, la Val d'Assa presenta un nucleo di incisioni di natura diversa dal Monte Luppia per contesto di scoperta, tecnica e grado di incertezza. Le incisioni furono esposte dalla grande piena del torrente Assa del 1966, che asportando terreno e vegetazione ne scoprì alcune e ne liberò parzialmente altre. La prima pubblicazione scientifica, firmata da Piero Leonardi con Giuseppe Rigoni e Aldo Allegranzi, risale al 1982 \citep{leonardi1982}. Il testo si apriva con una cautela significativa: si parlava di ``varie centinaia o forse addirittura migliaia (il rilevamento è appena iniziato)'' di incisioni, la maggior parte delle quali geometriche, eseguite a graffito, cioè per incisione leggera con uno strumento appuntito; una tecnica che produce segni meno profondi e quindi più difficili da datare.

Gli autori, riprendendo un'osservazione di Paolo Graziosi, concludevano che una parte consistente delle incisioni era probabilmente medievale, senza escludere che alcuni segni geometrici potessero risalire alla preistoria. La cautela non ha impedito la formazione di un'interpretazione più marcatamente preistorica del sito, sostenuta da Ausilio Priuli (1983), che ha mantenuto vivo un dibattito durato decenni, fino al convegno di Gallio e Canove del 1996 (atti 2001).

Un lavoro recente con tecniche digitali di rilievo tridimensionale e luce strutturata ha permesso una lettura più oggettiva delle superfici. La conclusione è che un'origine preistorica appare ``molto dubbia, o eventualmente limitata a segni geometrici che sono essenzialmente privi di tempo'', cioè non databili per loro stessa natura \citep{bettineschi2022}. Gran parte delle figure è di epoca medievale o moderna, comprese croci, raffigurazioni di chiese e diverse date leggibili, la più antica delle quali riporta l'anno 1506. Il sito è attualmente interdetto al pubblico per gravi dissesti del versante.

\subsection{I Colli Berici e la pianura vicentina}

I Colli Berici hanno restituito tracce di frequentazione risalenti al Paleolitico, in particolare nella Grotta di Broion e nella Grotta dei Ladroni, scavate da Luigi Capitan e poi da altri ricercatori nel corso del Novecento. La Grotta di Broion ha restituito industria musteriana associata a fauna tipica delle fasi fredde: rinoceronte lanoso, orso delle caverne, iena delle caverne \citep{bartolomei1980}. La frequentazione epigravettiana dei Berici è più frammentaria, con presenze attestate ma non ancora sistematicamente catalogate in anni recenti.

Il lago di Fimon, nella pianura vicentina, è noto per la presenza di un sito palafitticolo neolitico riportato alla luce da ricerche condotte tra gli anni sessanta e ottanta: ceramica della Cultura dei Vasi a Bocca Quadrata, resti vegetali carbonizzati e fauna domestica. È uno dei pochi siti lacustri del Veneto interno con una sequenza neolitica documentata con scavo.

% ============================================================
\section{Padova: gli Euganei, la pianura e il distretto di Este}
% ============================================================

La provincia di Padova presenta un quadro preistorico in cui montagna e pianura svolgono ruoli nettamente diversi, condizionati dalla geomorfologia: i Colli Euganei, gruppo di rilievi vulcanici isolati nella pianura veneta, hanno offerto risorse minerarie e punti di riferimento visivo eccezionali, mentre la pianura è stata abitata nelle sue fasi più recenti (bronzo, ferro) in modo intenso ma con una tafonomia complicata dalle alluvioni ripetute e dall'agricoltura.

\subsection{Paleolitico e Mesolitico: una presenza difficile da documentare}

Non risultano siti paleolitici con stratigrafia conservata nella provincia di Padova: i depositi pleistocenici della pianura sono stati ampiamente rimaneggiati dal reticolo fluviale e dalle bonifiche. La litica sporadica segnalata in letteratura per i versanti degli Euganei (Baone, Cinto Euganeo) non ha ancora prodotto contesti stratigrafici chiusi che permettano attribuzioni sicure.

La situazione mesolitica è analoga: l'assenza di siti con contesto è la regola, non l'eccezione. Un'eccezione parziale è il sito di Bosco dei Fontanassi a Piombino Dese, sul margine nordorientale della provincia, dove microliti geometrici riferibili al Mesolitico sono stati segnalati in raccolta di superficie in terreni agricoli; ma anche in questo caso la mancanza di contesto stratigrafico rende problematica qualsiasi interpretazione funzionale del sito.

\subsection{Neolitico: i villaggi della pianura}

Il Neolitico padovano è documentato principalmente dalla Cultura dei Vasi a Bocca Quadrata (VBQ), nella sua espressione pianura veneta, con villaggi attestati a Vigodarzere, Mejaniga e Campodarsego: insediamenti di pianura con strutture abitative, ceramica decorata, fauna domestica e industria litica levigata. Questi siti appartengono alla seconda metà del quarto millennio prima dell'anno zero e si collocano in un areale che dalla pianura padana si estende verso la costa adriatica.

Abano Terme, con le sue sorgenti termali calde, è considerata un possibile polo di attrazione già in età preistorica, ma le testimonianze dirette di frequentazione neolitica nell'area termale sono frammentarie, sovrapposte a fasi storiche successive e difficili da isolare senza scavi profondi che le infrastrutture moderne hanno reso impraticabili.

\subsection{Eneolitico: le risorse degli Euganei}

I Colli Euganei in età eneolitica (terzo millennio prima dell'anno zero) costituivano una fonte di materie prime di interesse sovraregionale: la trachite euganea era estratta e lavorata per la produzione di macine, pestelli e strumenti pesanti. Sono state segnalate presenze della Cultura del Vaso Campaniforme in contesti di altura sui Colli, ma il quadro resta poco sistematizzato, con materiali in gran parte da superficie o da scavi vecchi non documentati con standard moderni.

In questo contesto si inserisce una questione aperta: sulle superfici di trachite dei Colli Euganei sono state segnalate, in modo non sistematico, figure a coppella che potrebbero essere di età preistorica, per analogia con coppelle documentate su altri substrati vulcanici nell'arco alpino. Al momento in cui si scrive, non risulta una campagna di rilevamento sistematico che abbia affrontato questa questione con metodi comparabili a quelli usati per l'arte rupestre della Valcamonica o del Monte Luppia.

\subsection{Bronzo e Ferro: l'intensità della pianura}

L'età del bronzo padovana è documentata con una densità che in pianura non ha equivalenti nel Veneto montano. I dossi fluviali residui (paleoisole di sedimento emergenti dall'antica pianura alluvionale) hanno restituito materiali dell'età del bronzo a Padova città, Casalserugo, Legnaro e Este. Il Monte Cero a Baone (Colli Euganei) e la Rocca di Monselice sono siti di altura con frequentazione del bronzo medio e recente, con materiali difensivi e di stoccaggio. Il ritrovamento di piroghe monossili nel letto del Bacchiglione e del Brenta testimonia l'uso sistematico dei fiumi come vie di comunicazione già nel secondo millennio.

Il distretto di Este (Ateste) è il sito di riferimento per l'età del ferro del Veneto meridionale: necropoli con centinaia di tombe, situle decorate a sbalzo con scene narrative, e il santuario della dea Reitia con migliaia di offerte votive sono documentati con scavi sistematici dall'Ottocento \citep{chieco_bianchi1988}. Padova preromana, con il suo impianto urbano databile almeno al V secolo prima dell'anno zero, completa il quadro di una provincia con una densità di frequentazione storico-preistorica molto elevata nelle fasi più recenti.

Scoperte recentissime (2025) hanno portato alla luce una necropoli venetica di trecento tombe in via Campagnola a Padova, e un santuario venetico a Ponso: ulteriori conferme di un distretto densamente abitato durante l'età del ferro, con una produzione simbolica ricca di cui le situle sono il prodotto più noto, ma che non include arte rupestre documentata.

\subsection{Arte rupestre: assente come corpus}

Non risulta arte rupestre documentata nella provincia di Padova. Le coppelle sugli Euganei, già citate, rappresentano un punto interrogativo aperto piuttosto che un corpus attestato. Questa assenza, come si è detto a proposito della tafonomia nel capitolo precedente, non equivale a una prova di assenza: equivale a documentare lo stato della ricerca, che su questo aspetto specifico non ha ancora prodotto un inventario sistematico.

% ============================================================
\section{Rovigo: il Polesine e la sua preistoria sommersa}
% ============================================================

La provincia di Rovigo presenta una situazione preistorica strutturalmente diversa da tutte le altre province venete, condizionata in modo determinante dalla sua geomorfologia: il Polesine è una pianura alluvionale depressa, compresa tra Po e Adige, che nel corso della preistoria era un paesaggio di paludi, dossi fluviali e lagune, in continua trasformazione a causa delle variazioni del livello marino e delle divagazioni fluviali. Questa geomorfologia ha una conseguenza diretta per l'archeologia: i siti più antichi, quando esistono, sono sepolti a profondità variabili sotto i sedimenti alluvionali e non sono accessibili con ricognizioni di superficie.

\subsection{Paleolitico e Mesolitico: assenza strutturale}

Non risultano siti paleolitici o mesolitici documentati con contesto stratigrafico nella provincia di Rovigo. Questa non è un'eccezione: è la norma attesa per un territorio che, nelle fasi più antiche del Pleistocene superiore, era il delta di un sistema fluviale maggiore, e nelle fasi del Mesolitico era probabilmente ancora in gran parte sommerso o lagunare, con la linea di costa a una posizione notevolmente più interna rispetto all'attuale. I depositi che potrebbero contenere testimonianze di frequentazione paleolitica si trovano a profondità non indagabili con archeologia convenzionale.

\subsection{Neolitico ed Eneolitico: la prima presenza documentata}

Le prime testimonianze di frequentazione documentata nel Polesine risalgono al Neolitico, con ceramica raccolta in superficie sui dossi fluviali residui (le uniche quote positive del territorio). La rivista \emph{Padusa}, pubblicata dal Centro Polesano di Studi Storici Archeologici ed Etnografici di Rovigo, costituisce il principale strumento di censimento della letteratura locale, con segnalazioni che si estendono dal Neolitico fino all'età medievale \citep{padusa}. La qualità della documentazione è però variabile: molte segnalazioni derivano da rinvenimenti fortuiti in lavori agricoli o canalizzazioni, senza documentazione stratigrafica.

\subsection{Età del bronzo: Frattesina e il distretto metallurigico}

L'età del bronzo recente e finale (XII-X secolo prima dell'anno zero) ha lasciato nel Polesine tracce di un'intensità eccezionale nel panorama dell'Italia settentrionale. Il sito di Frattesina di Fratta Polesine, scavato sistematicamente dall'Università di Ferrara, copre una superficie stimata in oltre venticinque ettari e ha restituito più di mille tombe, officine per la lavorazione dell'avorio, del vetro e dell'ambra, e la prima fornace vetraria documentata in Europa occidentale \citep{bietti_sestieri1996}. Il sito di Campestrin a Grignano Polesine ha restituito oggetti in ambra baltica, prova di contatti commerciali con un'area distante migliaia di chilometri. Il sito di Canàr a Castelnovo Bariano era un insediamento su palafitta lungo il Po, con tipologie ceramiche di contatto tra il mondo dell'Italia settentrionale e quello transalpino. Villamarzana ha restituito una necropoli con cremazioni della stessa fase.

\subsection{Età del ferro: San Basilio di Ariano e Adria}

San Basilio di Ariano Polesine è un emporio dell'età del ferro (VII-V secolo prima dell'anno zero) con testimonianze di ceramica greca, etrusca e venetica: un punto nodale nella rete commerciale adriatica \citep{bonomi1991}. Adria, che ha dato il nome all'Adriatico, è una città portuale con una lunga frequentazione preistorica e storica; i materiali più antichi dal suo territorio risalgono all'età del bronzo, ma è in età arcaica e classica che il sito raggiunge la massima importanza come porto di transito.

\subsection{Arte rupestre: impossibile per ragioni geomorfologiche}

Nel Polesine l'arte rupestre è strutturalmente impossibile: mancano rocce affioranti. Il substrato è ovunque coperto da sedimenti alluvionali profondi, e nessuna superficie rocciosa è accessibile in modo naturale in tutta la provincia. Questa non è una lacuna della ricerca: è una conseguenza della geologia. La ricchissima preistoria del Polesine si esprime attraverso insediamenti, necropoli e manufatti mobili, non attraverso arte su roccia.

% ============================================================
\section{Venezia: la laguna come risultato, non come sfondo}
% ============================================================

La provincia di Venezia richiede una premessa geomorfologica senza la quale i dati archeologici non si capiscono. La laguna che oggi identifica il territorio non è un paesaggio immutabile: è il risultato finale di un processo durato diecimila anni. Nel Tardoglaciale, attorno ai diciottomila anni prima dell'anno zero, il livello del mare era circa centoventi metri più basso dell'attuale e tutta l'area dell'Alto Adriatico era una vasta pianura emersa, percorsa dal Brenta, dall'Adige e da altri fiumi che defluivano verso un litorale molto più a est. La trasgressione Flandriana, tra i diecimila e i cinquemilacinquecento anni prima dell'anno zero, ha progressivamente sommerso quella pianura. Circa cinquemilacinquecento anni fa la linea di costa si trovava trenta chilometri più a ovest dell'attuale, e i dossi fluviali del Brenta (Mirano, Spinea, Noale, Martellago) erano già gli unici rilievi sabbiosi emergenti nell'area perilagunare. La laguna si è strutturata nella sua forma attuale solo tra quattromila e tremila anni prima dell'anno zero, in età eneolitica e del bronzo.

Questa cronologia ha una conseguenza diretta per l'archeologia: i siti più antichi, se sono stati occupati nelle fasi in cui il territorio era ancora pianura emersa, si trovano oggi sepolti sotto metri di sedimento lagunare o marino, non accessibili con metodi di ricerca convenzionali.

\subsection{Paleolitico e Mesolitico: assenza strutturale e tracce di superficie}

Non risultano siti paleolitici documentati nella provincia di Venezia; la spiegazione è la stessa del Polesine: i depositi pleistocenici sono sepolti sotto le alluvioni successive e non affiorano. Per il Mesolitico la situazione è parzialmente diversa: prospezioni di superficie a partire dal 1975 nella terraferma perilagunare, nell'area di Mestre, hanno individuato materiale litico in selce riferibile al Mesolitico finale (attorno ai seimilacinquecento anni prima dell'anno zero). Il limite di questi rinvenimenti è lo stesso che si incontra altrove nella pianura veneta: si tratta di materiale raccolto in superficie in terreni agricoli, senza contesto stratigrafico chiuso, e la sua interpretazione funzionale rimane problematica.

\subsection{Neolitico: Concordia Sagittaria}

Il sito di Concordia Sagittaria, in particolare la località Loncon, ha restituito nel 2012 un insediamento tardo-neolitico databile alla metà del quarto millennio prima dell'anno zero, con fosse domestiche, vasche e canali, seguito da una fase del bronzo antico. La pubblicazione in \emph{Notizie di Archeologia del Veneto} è uno dei pochi casi in cui un sito della provincia veneziana è documentato con scavo stratigrafico e cronologia assoluta. Per il resto del territorio neolitico, i dossi del Brenta (Mirano-Spinea in particolare) sono morfologicamente favorevoli ma non ancora indagati sistematicamente.

\subsection{Età del bronzo: le piroghe del Brenta meridionale}

Un ritrovamento del 2014 nel canale del Cornio a Campagna Lupia ha portato alla luce scafi monossili (piroghe in quercia) nell'antico alveo del Brenta meridionale, pubblicati sulla rivista del Museo Civico di Storia Naturale di Venezia \citep{asta2014}. Il contesto, con resti di castoro che indicano un ambiente forestale-ripario oggi scomparso, permette di fissare la frequentazione del corso d'acqua in un'epoca in cui la laguna non esisteva ancora nella sua forma attuale e i fiumi defluivano liberamente verso il mare.

\subsection{Età del ferro: Altino e i Veneti antichi}

Il sito più importante della provincia per la preistoria recente è Altino (Quarto d'Altino), che ha restituito una sequenza occupazionale continua dal bronzo recente (XIII-XII secolo prima dell'anno zero) fino alla romanizzazione. Tra il IX e l'VIII secolo prima dell'anno zero era già uno scalo fluviale venetico; tra il VII e il VI secolo aveva assunto una fisionomia urbana con un insediamento di circa quaranta ettari. Il Museo Nazionale di Altino, con la sua collezione di materiali venetici, è il riferimento principale per chi voglia capire cosa significa, in termini di produzione simbolica, la fase che precede immediatamente la romanizzazione in quest'area. Concordia Sagittaria presenta fasi venetiche analoghe, più a est.

\subsection{Arte rupestre: assente per ragioni strutturali}

Come nel Polesine, nell'intera provincia di Venezia non esistono affioramenti rocciosi naturali. L'arte rupestre è strutturalmente impossibile, non una lacuna della ricerca. La preistoria veneziana si esprime attraverso insediamenti, piroghe, ceramica e necropoli: non attraverso segni sulla roccia.

% ============================================================
\section{Treviso: la Marca e le Prealpi}
% ============================================================

% SEGNAPOSTO

\emph{Sezione in preparazione. La provincia di Treviso è inclusa nella pianura trevigiana e nelle Prealpi Trevigiane (Cansiglio, Cesen, Nevegal), con frequentazione mesolitica documentata nell'area del Cansiglio.}

% ============================================================
\section{Belluno: le Dolomiti e il Cadore}
% ============================================================

% SEGNAPOSTO: la ricerca sistematica del Cadore merita trattazione estesa

\emph{Sezione in preparazione. La provincia di Belluno, con il Cadore e le Dolomiti, è trattata parzialmente nel capitolo~\ref{cap:dove_guardare} a proposito dei casi di successo recente: l'Uomo di Mondeval, i siti mesolitici di Val Visdende e i ritrovamenti di Pian dei Buoi. Questa sezione integrerà il quadro con i siti del Bellunese orientale (Alpago, Zoldo, Agordino) e le grotte del massiccio delle Vette.}

% ============================================================
\section{Il Veneto geografico esteso: Garda trentino e lombardo, Primiero, Folgaria}
% ============================================================

% SEGNAPOSTO

\emph{Sezione in preparazione. La sponda trentina (Riva del Garda, Arco, Torbole) e la sponda lombarda settentrionale del Garda hanno restituito arte rupestre dell'età del bronzo in continuità con il complesso del Monte Luppia. Primiero e Folgaria sono inclusi per la continuità geografica con i valichi alpini frequentati in preistoria, non per ragioni amministrative.}

% ============================================================
\section{Un confronto numerico}
% ============================================================

I siti e i complessi descritti in questo capitolo acquistano significato pieno solo se messi accanto ai numeri delle aree più studiate d'Italia e d'Europa.

\begin{table}[htbp]
\centering
\small
\begin{tabular}{@{}p{2.6cm}p{1.6cm}p{1.8cm}p{1.6cm}p{2.4cm}@{}}
\toprule
\textbf{Sito} & \textbf{Rocce} & \textbf{Figure} & \textbf{Inizio ricerca} & \textbf{Cronologia proposta} \\
\midrule
Valcamonica (BS) & $\sim$2.400 & $>$140.000 & anni 1930 (sistematica dal 1964) & dal Neolitico al periodo storico \\
Monte Bego (Francia) & migliaia & decine di migliaia & primi '900 & età del bronzo \\
Garda/Monte Luppia (VR) & $>$250 & $\geq$3.000 & 1964 & età del bronzo -- epoca moderna \\
Val d'Assa (VI) & non quantificato & centinaia/migliaia (stima 1982) & 1966/1982 & prevalentemente medievale/moderna \\
Alto Lario (CO) & $\sim$300 & non quantificato & anni 1970 & età del ferro \\
\bottomrule
\end{tabular}
\caption{Confronto tra i principali complessi di arte rupestre alpina citati in questo libro. I dati sul Garda derivano da Gaggia (1982, 2002); quelli sulla Val d'Assa da Leonardi, Rigoni e Allegranzi (1982) e dalla revisione di Bettineschi e collaboratori (2022); quelli sulla Valcamonica dalla documentazione presentata all'UNESCO nel 1979 e da lavori successivi.}
\label{tab:confronto_siti}
\end{table}

Alcuni aspetti di questa tabella meritano un commento esplicito. La colonna delle figure per la Val d'Assa non è comparabile con le altre: mentre per Garda e Valcamonica il numero deriva da un catalogo sistematico di figure attribuite a un'epoca, per la Val d'Assa è una stima complessiva di tutti i segni visibili, indipendentemente dalla loro epoca; la maggior parte di quei segni è oggi ritenuta medievale o moderna. La colonna dell'inizio della ricerca mostra un divario che nessun'altra colonna può spiegare: la Valcamonica ha alle spalle quasi un secolo di attenzione scientifica; il Garda ha sessant'anni, iniziati da un'intuizione personale innescata dal contatto con la Valcamonica; la Val d'Assa ha un rilevamento iniziato per caso, in seguito a un'alluvione, e mai completato con la stessa sistematicità.

% ============================================================
\section{Cosa questo quadro dimostra, e cosa no}
% ============================================================

Va detto con la massima chiarezza possibile cosa questo capitolo non dimostra. Non dimostra che il Veneto fosse, in antichità, meno frequentato della Lombardia o delle Alpi occidentali: potrebbe esserlo stato, ma il dato numerico da solo non lo prova. Quello che il confronto mostra con certezza è una differenza nella storia della ricerca, misurabile e non solo impressionistica: alcune aree (Lessinia per il Paleolitico, distretto di Este per l'età del ferro, Polesine per l'età del bronzo) hanno ricevuto un'attenzione sistematica di livello europeo, e i risultati sono stati proporzionali all'investimento. Altre aree, in particolare la fascia alpina e prealpina centrale, hanno ricevuto ricerche discontinue o fortuite, e il quadro che ne emerge riflette questa discontinuità più che la realtà preistorica.

La Valcamonica ha alle spalle una scuola di studi continuativa, nata negli anni sessanta del Novecento attorno al lavoro di Emmanuel Anati e mai interrotta, con generazioni di ricercatori dedicati esclusivamente a quel territorio. Il Veneto non ha mai avuto un equivalente per l'arte rupestre. Le sue scoperte più importanti in questo campo, come si è visto, sono nate da episodi in larga parte fortuiti: una visita che accende un'intuizione, una piena che scopre una parete. Questa differenza, da sola, non spiega necessariamente tutto il divario numerico osservato. Ma è sufficiente per rendere prudente qualunque conclusione che dia per scontata l'assenza di ciò che, semplicemente, non è ancora stato cercato con la stessa intensità.
